\section{Lecture 15: Toward Analytic Stacks --- Motivation from Algebraic Geometry}
\label{lec:15}

Having built the theory of analytic rings and worked through a range of examples
(solid, liquid, gaseous; \cref{lec:8,lec:12,lec:13,lec:14}), Clausen turns to the
next stage of the course: using analytic rings to build \emph{geometric}
objects, the \emph{analytic stacks}. This lecture is entirely motivational --- the
precise definition of an analytic stack is deferred --- and proceeds by analogy
with classical algebraic geometry. Everything is kept deliberately loose; several
statements are heuristic, and two live corrections (recorded below) reshaped the
account of analytification and of GAGA over $\mathbb{Q}_p$ as the lecture went
on. The plan:
\begin{itemize}
\item recall the modern (functor-of-points) paradigm for schemes and stacks, and
the question it answers: how to single out a good class of geometric objects
inside sheaves on affines;
\item survey why one wants \emph{stacks} --- and objects even more general than
Artin stacks --- and arrive at the (perhaps surprising) answer that the right
ambient category is \emph{all} sheaves on affines, modulo set theory;
\item transport this to analytic geometry: a Grothendieck topology on
$(\operatorname{AnRing})^{\mathrm{op}}$, the desiderata one wants the resulting
stacks to satisfy, the mechanism of \emph{analytification}, and GAGA read as an
isomorphism of stacks;
\item address the set-theoretic technicalities via accessible presheaves.
\end{itemize}

Throughout, ``ring'' means commutative ring, and we work over a loosely-fixed base
(usually $\mathbb{Z}$, occasionally a field). We write $\operatorname{Aff} =
(\operatorname{CRing})^{\mathrm{op}}$ for affine schemes.

\subsection{The paradigm of algebraic geometry}

The paradigm is that commutative rings describe basic geometric objects, the
\emph{affine schemes}, via the (anti-)equivalence
\begin{equation}
\label{eq:15:spec}
R\ \longmapsto\ \operatorname{Spec} R,\qquad \operatorname{CRing}\ \simeq\ \operatorname{Aff}^{\mathrm{op}},
\end{equation}
and that one wants to \emph{glue} affine schemes together to make more general
geometric objects.

\begin{remark}[Two aspects of gluing]
\label{rmk:15:two-aspects}
Gluing has two aspects: \emph{(1) which gluings are allowed}, and \emph{(2) how to
identify the results of two different gluings} --- more generally, what the maps
between formal gluings are, i.e.\ what the category is. Aspect~(2) is best confronted
first: the projective line $\mathbb{P}^1 = \mathbb{A}^1_+ \cup_{\Gm} \mathbb{A}^1_-$
(glue two copies of $\mathbb{A}^1$ along $\Gm$ by inversion) must acquire symmetries
such as $\mathrm{PGL}_2$ that do \emph{not} respect the chosen decomposition, so the
category must be told about maps invisible to any one cover.
\end{remark}

The cleanest way to handle~(2) is to \emph{specify an ambient category} containing
$\operatorname{Aff} = (\operatorname{CRing})^{\mathrm{op}}$, and single out a full
subcategory of ``allowable gluings of affine schemes'' inside it.

\begin{remark}[Classical vs.\ modern ambient categories]
\label{rmk:15:classical-modern}
There are two approaches.
\begin{itemize}
\item \emph{Classical (schemes).} Take the ambient category to be locally ringed
spaces; $R\mapsto \operatorname{Spec} R$ is such, and a scheme is a locally ringed
space locally isomorphic to some $\operatorname{Spec} R$. The allowable gluings of
affine opens are gluings along open subsets in the locally-ringed-space sense.
\item \emph{Modern.} Do not try to be clever about the ambient category; build it
formally from $\operatorname{CRing}$. The universal category in which one can glue
(i.e.\ freely adjoin colimits) is the category of presheaves
$\operatorname{PSh}(\operatorname{Aff})$: giving a colimit-preserving functor out of
it is the same as giving an ordinary functor from $\operatorname{Aff}$. (There is a
set-theoretic caution here --- $\operatorname{Aff}$ is not small --- taken up in
\cref{subsec:15:set-theory}.) One has \emph{formally} allowed gluing but not yet said
how to identify the results of two gluings; that is done by passing to sheaves
$\operatorname{Sh}(\operatorname{Aff})$ for a Grothendieck topology. More covers in
the topology means more maps get made (like the $\mathrm{PGL}_2$-symmetries of
$\mathbb{P}^1$); too many covers destroys information, so the choice of topology is
delicate.
\end{itemize}
\end{remark}

\begin{example}[Schemes, modern-style]
\label{ex:15:schemes}
For the Zariski topology one has
\begin{equation}
\label{eq:15:schemes-in-sheaves}
\{\text{schemes}\}\ \subseteq\ \operatorname{Sh}^{\mathrm{Zar}}(\operatorname{Aff})
\end{equation}
as the full subcategory of Zariski sheaves that are \emph{locally representable} in
the sense of open covers; equivalently, one defines open inclusions in
$\operatorname{Sh}^{\mathrm{Zar}}(\operatorname{Aff})$ by reduction (via pullback) to
the affine case. This is the perspective we will take for analytic stacks. But
schemes are too restrictive a model: more general geometric objects come up and are
genuinely relevant.
\end{example}

\subsection{Why stacks?}
\label{subsec:15:why-stacks}

Often, moduli spaces in algebraic geometry have the form $X/G$ with $X$ a variety
and $G$ a group variety (or scheme) acting on $X$.

\begin{example}[Moduli of elliptic curves]
\label{ex:15:m-ell}
Inverting $6$ for simplicity, the moduli of elliptic curves is
\begin{equation}
\label{eq:15:m-ell}
\mathcal{M}_{\mathrm{ell},\,\mathbb{Z}[1/6]}\ =\ \operatorname{Spec}\bigl(\mathbb{Z}[1/6][A,B][\Delta^{-1}]\bigr)\big/\Gm ,
\end{equation}
parametrizing the elliptic curve in Weierstrass form $y^2 = x^3 + Ax + B$ (invert
the discriminant $\Delta$); the only isomorphisms between two such curves are scalar
substitutions with fixed weights on $x,y$, giving the $\Gm$-action. The same stack
has other presentations, e.g.\ after adding level structure,
$\mathcal{M}_{\mathrm{ell},N}/\mathrm{GL}_2(\mathbb{Z}/N\mathbb{Z})$ (with
$\mathcal{M}_{\mathrm{ell},N}$ a variety for $N$ large); the chosen Grothendieck
topology must identify the two, so one at least allows \'etale covers --- the classic
choice.
\end{example}

\begin{remark}[The naive quotient is bad]
\label{rmk:15:bad-quotient}
The quotients of \cref{ex:15:m-ell} exist already in schemes, giving
$\mathbb{A}^1_{\mathbb{Z}[1/6]}$, the ``$j$-line'' (via the $j$-invariant). But this
is \emph{not a good quotient}. On $\mathcal{M}_{\mathrm{ell}}$ there is a natural
line bundle $\omega = \operatorname{Lie}(E_{\mathrm{univ}})^{*}$ (the cotangent space
at the origin of the universal elliptic curve), which is $\Gm$-equivariant upstairs
but does not descend to $\mathbb{A}^1$; and the universal elliptic curve itself
cannot be defined over $\mathbb{A}^1$. The problem is that the action is \emph{not
free} (some elliptic curves have extra automorphisms), so the scheme-theoretic
quotient collapses too much and forgets information.
\end{remark}

\begin{remark}[Groupoid quotients and Artin stacks]
\label{rmk:15:groupoid-quotient}
The solution is to take the quotient in a more refined ($2$-)category, that of
(sheaves of) \emph{groupoids}. The groupoid quotient $X/G$ pretends every action is
free: its objects are points of $X$ and its morphisms are elements of $G$, so a
fixed point acquires extra automorphisms, and the map $X\to X/G^{\mathrm{gpd}}$
always behaves like the total space of a $G$-bundle --- fibres (taken as
$2$-categorical fibre products) are all isomorphic to $G$:
\begin{equation}
\label{eq:15:groupoid-fibre}
\begin{tikzcd}
G \arrow[r] \arrow[d] & X \arrow[d] \\
\ast \arrow[r] & X/G^{\mathrm{gpd}}
\end{tikzcd}
\end{equation}
Line bundles on $X/G^{\mathrm{gpd}}$ are then exactly $G$-equivariant line bundles on
$X$. This leads to \emph{Deligne--Mumford} and, more generally, \emph{Artin} stacks:
full subcategories of \'etale sheaves on $\operatorname{Aff}$ that are, roughly,
locally --- in the sense of admitting a \emph{smooth} cover by affine schemes ---
of the form $X/G$ for a smooth group scheme $G$ acting on an affine $X$. This is
excellent for moduli theory; the theoretical underpinning is Artin's
representability theorem \cite{Artin_1974}.
\end{remark}

\begin{remark}[The constraint: smooth covers force finite type]
\label{rmk:15:smooth-forces-ft}
Artin stacks are still constrained by the need for a smooth cover, which in
particular forces \emph{finite presentation}. An infinite-dimensional (non-finite-type)
group scheme acting on something need not admit a quotient in this world.
\end{remark}

There are many geometrically relevant $X\in \operatorname{Sh}^{\text{\'et}}(\operatorname{Aff})$
that are \emph{not} Artin stacks (nor schemes).

\begin{example}[Formal schemes]
\label{ex:15:formal-schemes}
For $R$ noetherian and $I\subseteq R$ an ideal, one wants the formal spectrum. Huber's
theory (\cref{lec:8}) and Grothendieck's theory both view it as a topological ring
(localizing along different subsets); but another viewpoint is simply the union of
the infinitesimal neighbourhoods,
\begin{equation}
\label{eq:15:spf}
\operatorname{Spf}(R^{\wedge}_I)\ =\ \bigcup_{n}\ \operatorname{Spec}(R/I^{n}),
\end{equation}
so that giving finite-type data on the formal scheme --- e.g.\ a vector bundle ---
is giving a compatible collection over all the thickenings:
\begin{equation}
\label{eq:15:vect-formal}
\operatorname{Vect}(R^{\wedge}_I)\ =\ \varprojlim_{n}\ \operatorname{Vect}(R/I^{n}).
\end{equation}
This is a very different gluing from those of schemes or Artin stacks: no smooth
covers are in sight, only a union of infinitesimal thickenings.
\end{example}

\begin{example}[Moduli of formal groups]
\label{ex:15:formal-groups}
For some moduli problems one genuinely quotients by an infinite-dimensional group.
The example dear to homotopy theorists (and number theorists) is the moduli of
$1$-dimensional formal groups $\mathcal{M}_{\mathrm{FG}}$ (with the Lubin--Tate
spaces as covers): the group to mod out by is coordinate changes on a
$1$-dimensional formal scheme, with infinitely many coefficients, hence
infinite-dimensional. This does not fit the Artin framework and suggests moving to a
finer topology, e.g.\ \emph{fpqc} rather than \'etale, to accommodate infinite-type
covers.
\end{example}

\begin{remark}[Simpson's slogan]
\label{rmk:15:simpson}
There is a remarkable class of examples originating with Carlos Simpson, which one
might crudely summarize --- Clausen stresses this is his own too-strong-to-be-literally-true
heuristic, not a claim of Simpson's --- as: \emph{every linear-algebra category is}
$\operatorname{QCoh}$ \emph{of some (suitably generalized) stack}. Simpson's work
produced many instances of natural linear-algebra categories realized as
quasi-coherent sheaves on a stack.
\end{remark}

\begin{example}[Some Simpson-type stacks]
\label{ex:15:simpson-stacks}
Let $k$ be a field.
\begin{enumerate}
\item For an algebraic group $G/k$, the classifying stack $BG = \ast/G^{\mathrm{gpd}}$
has
\begin{equation}
\label{eq:15:BG}
\operatorname{QCoh}(BG)\ =\ \operatorname{Rep}_G(k\text{-}\mathrm{Vect}),
\end{equation}
$G$-equivariant $k$-vector spaces, i.e.\ representations of $G$.
\item For $\mathbb{A}^1/\Gm$ (scalar multiplication), $\Gm\subseteq\mathbb{A}^1$ is a
$\Gm$-invariant open on which one takes $\Gm/\Gm = \ast$, while the closed complement
is $0/\Gm = B\Gm$; so this stack has a point as an open substack and $B\Gm$ as its
closed complement. Then
\begin{equation}
\label{eq:15:filtered}
\operatorname{QCoh}^{\mathrm{flat}}(\mathbb{A}^1/\Gm)\ =\ \{\,\text{$k$-vector spaces with a $\mathbb{Z}$-indexed filtration}\,\}.
\end{equation}
\end{enumerate}
\end{example}

\begin{example}[The de Rham stack]
\label{ex:15:de-rham}
Let $k$ have characteristic $0$ and $X/k$ a smooth variety. Simpson's de Rham stack
$X^{\mathrm{dR}}$ is presented by quotienting $X$ not by a group action but by the
equivalence relation identifying infinitesimally close points: one takes the formal
completion of $X\times_k X$ along the diagonal,
% board 18 @ 00:31:29 --- X^dR presented by the formal completion of the diagonal in X x_k X
\begin{equation}
\label{eq:15:de-rham-presentation}
\begin{tikzcd}
(X\times_k X)^{\wedge}_{X} \arrow[r, shift left=0.6ex] \arrow[r, shift right=0.6ex] & X \arrow[r] & X^{\mathrm{dR}},
\end{tikzcd}
\end{equation}
the completion taken in the sense of \eqref{eq:15:spf}. Then
\begin{equation}
\label{eq:15:vect-de-rham}
\operatorname{Vect}(X^{\mathrm{dR}})\ =\ \{\,\text{vector bundles on $X$ with flat connection}\,\},
\end{equation}
Grothendieck's interpretation of a flat connection as descent data along
infinitesimally close points, and the cohomology of the structure sheaf recovers de
Rham cohomology,
\begin{equation}
\label{eq:15:de-rham-cohomology}
R\Gamma(X^{\mathrm{dR}}, \mathcal{O})\ =\ R\Gamma_{\mathrm{dR}}(X/k).
\end{equation}
This is far from an Artin stack: one is modding out by a formal scheme giving an
equivalence relation.
\end{example}

\begin{remark}[Prismatic and related stacks]
\label{rmk:15:blf}
More recently, Bhatt--Lurie and Drinfeld define stacks whose $\operatorname{QCoh}$
captures coefficient systems for various $p$-adic cohomology theories in
characteristic $p$, mixed characteristic, or in the limit over $\mathbb{Z}_p$
(\cite{drinfeld2024prismatization,bhatt2022absoluteprismaticcohomology}). For de Rham
cohomology in characteristic $p$ there is again a stack --- using the divided-power
envelope of the diagonal in place of the formal neighbourhood
\eqref{eq:15:de-rham-presentation}. Prismatic, de Rham, and crystalline cohomology,
and the comparison theorems between them, then acquire a geometric explanation: an a
priori linear-algebraic comparison of cohomology theories is promoted to an
isomorphism of stacks, from which the comparison follows by passing to quasi-coherent
sheaves and functoriality.
\end{remark}

\begin{example}[Betti stacks and condensed sets]
\label{ex:15:betti}
Betti cohomology --- singular/sheaf cohomology with constant coefficients on the
analytic topological space underlying a complex variety --- also has a stacky
avatar. For $S$ a compact Hausdorff space, choose a surjection $T_0\twoheadrightarrow S$
from a profinite set; then $T_1 := T_0\times_S T_0$ is again profinite (a closed
subset of $T_0\times T_0$), so $S$ is the quotient of an equivalence relation in
profinite sets. Applying $\operatorname{Spec}$ of the ring of continuous
$\mathbb{Z}$-valued (or $k$-valued) functions gives a groupoid in affine schemes,
whose quotient in fpqc sheaves is the Betti stack:
% board 20/21 @ 00:38:07 --- Betti stack as fpqc quotient of the profinite presentation of S
\begin{equation}
\label{eq:15:betti-presentation}
\begin{tikzcd}
\operatorname{Spec} C(T_1,\mathbb{Z}) \arrow[r, shift left=0.6ex] \arrow[r, shift right=0.6ex] & \operatorname{Spec} C(T_0,\mathbb{Z}) \arrow[r] & S^{\mathrm{Betti}}.
\end{tikzcd}
\end{equation}
Then quasi-coherent sheaves on $S^{\mathrm{Betti}}$ are ordinary sheaves of
$\mathbb{Z}$-modules (abelian groups) on the topological space $S$, and coherent
cohomology on $S^{\mathrm{Betti}}$ is topological cohomology of $S$. More generally a
condensed set can be viewed as a stack, by mapping its profinite presentation into
the world of stacks via this procedure:
\begin{equation}
\label{eq:15:condensed-to-stacks}
\{\text{condensed sets}\}\ \longrightarrow\ \{\text{stacks}\}.
\end{equation}
\end{example}

\begin{remark}[Why the descent works]
\label{rmk:15:profinite-descent}
The checks for \eqref{eq:15:betti-presentation}: a surjection of profinite sets
induces a faithfully flat map on continuous functions, and fibre products of
profinite sets correspond to relative tensor products of the function rings. Both
reduce, via filtered colimits, to finite sets, where $C(-,\mathbb{Z})$ is a finite
product of copies of $\mathbb{Z}$ and everything is purely combinatorial. In fact any
map of profinite sets is flat (surjective $\Rightarrow$ faithfully flat), and
$\operatorname{Spec} C(T_0,\mathbb{Z})$ is $\underline{T_0}$, the inverse limit of the
constant schemes on the finite quotients.
\end{remark}

Faced with this baffling array of stacks --- many bearing no resemblance to Artin
stacks, yet wanted as convenient encodings of linear-algebraic and geometric
phenomena --- one asks:

\begin{remark}[The question and its answer]
\label{rmk:15:answer}
\emph{Question.} How to define a reasonable full subcategory of
$\operatorname{Sh}^{\mathrm{fpqc}}(\operatorname{Aff})$ containing all of the above
examples?

\emph{Answer.} $\operatorname{Sh}^{\mathrm{fpqc}}(\operatorname{Aff})$ itself --- the
whole category (modulo set theory, \cref{subsec:15:set-theory}). It is not clear that
there is any other answer. There is content in this answer, but it introduces
nothing not already present in the question: one takes the full subcategory that is
the whole thing.
\end{remark}

\subsection{Analytic stacks: the plan and the desiderata}

\begin{remark}[The plan]
\label{rmk:15:plan}
We take the same approach for analytic stacks: define a Grothendieck topology on the
affine analytic stacks $(\operatorname{AnRing})^{\mathrm{op}}$ --- the opposite of
analytic rings --- and take sheaves with respect to it (again modulo set theory). We
will not fix the precise topology today; instead we describe the phenomena the
resulting theory should capture.
\end{remark}

\begin{remark}[Desiderata]
\label{rmk:15:desiderata}
The examples we want (each producing an analytic stack):
\begin{enumerate}
\item \emph{Adic spaces} (Huber). The building blocks --- Huber pairs $(R,R^+)$ ---
already give analytic rings (\cref{lec:8,lec:9}), and the localization of
$D(R)$ over the Huber spectrum (\cref{lec:9,lec:10}) is the input to gluing along
rational open subsets. We want that gluing to yield an analytic space.
\item \emph{Complex analytic spaces}, over a gaseous or $p$-liquid version of
$\mathbb{C}$, with gluing along complex-analytic open subsets allowed.
\item \emph{Algebraic stacks}: analytic geometry should generalize schemes. Over the
uncompleted base $\mathbb{Z}$ (the trivial/maximal structure, with $D = D(\ur R)$ for
all modules, cf.\ \cref{rmk:1:trivial-structure}) one can do ordinary algebraic
geometry; a slight modification of \cref{rmk:15:answer} gives algebraic stacks as
analytic stacks. Universally over any analytic ring, an algebraic object yields an
analytic object by base change.
\item \emph{Banach rings} $\rightsquigarrow$ an analytic stack matching Berkovich's
theory. Interestingly, the stack attached to a general Banach ring is \emph{not}
affinoid but a genuine stack: e.g.\ $\mathbb{Z}$ with its usual archimedean norm goes
to a stack corresponding to the Berkovich spectrum $\MBerk(\mathbb{Z})$ (a tree,
cf.\ \cref{prop:12:berkovich-Z}), not an affinoid.
\item \emph{Coefficient systems}: analytic stacks whose $\operatorname{QCoh}$
captures various cohomology-coefficient systems (as in \cref{rmk:15:blf}).
\item \emph{The Tate curve} and its uniformization over the gaseous base ring
(\cref{def:14:tate-construction}), together with machinery to prove it algebraic ---
a Riemann--Roch theorem (giving a projective embedding from the identity section)
and a GAGA theorem.
\end{enumerate}
\end{remark}

\subsection{Analytification: importing Huber's theory of pairs}

We spend some time on how one item --- GAGA --- can be explained as an isomorphism
of stacks, since it illustrates the mechanism.

\begin{remark}[Recollection: the solid affine line]
\label{rmk:15:recall-solid-disk}
Adic geometry in the solid context (\cref{lec:7}) got started once we saw that there
is a nice ``subset'' of the affine line $\mathbb{A}^1$ over $\Zsolid$, namely the
closed unit disk (which \cref{rmk:7:verdict} taught us to think of as \emph{open}).
Algebraically it came from an \emph{idempotent} algebra --- the power series ring
$\mathbb{Z}[\![T]\!] = P^{\solid}$, regarded in $\mathbb{Z}[T]$-modules in solid
abelian groups (the analytic ring controlling $\mathbb{A}^1_{\Zsolid}$,
\cref{lem:7:idempotent,ex:7:closed-disk-Qp}). Moving it to infinity by $T\mapsto
T^{-1}$ produced the closed unit disk, which lets one tie into the $(R,R^+)$-theory:
one asks that the elements of $R^+$ actually land in the closed unit disk, not merely
map to $\mathbb{A}^1$.
\end{remark}

\begin{remark}[The gaseous analogue: germs at the origin]
\label{rmk:15:germs}
Working instead over the gaseous ring $\mathbb{Z}[\widehat{q}^{\,\pm1}]^{\mathrm{gas}}$
(so $q$ is inverted; cf.\ \cref{thm:14:gaseous-base}), one can again define a
``subset'' of $\mathbb{A}^1_{\mathbb{Z}[\widehat q^{\pm1}]^{\mathrm{gas}}}$
corresponding to an idempotent algebra --- morally, the functions convergent in
\emph{some unspecified} disk around the origin, i.e.\ germs of functions at $0$. It is
a filtered colimit
\begin{equation}
\label{eq:15:germ-colimit}
\varinjlim\bigl(\,P \xrightarrow{\ q\,\cdot\ } P \xrightarrow{\ q\,\cdot\ } \cdots\,\bigr)
\end{equation}
of the basic object $P$ under multiplication by $q$. In contrast to the solid case,
where $P$ ($=\mathbb{Z}[\![T]\!]$ once solidified) was already idempotent, here $P$ is
\emph{not} idempotent as an object of $\mathbb{Z}[T]$-modules in the gaseous category;
but taking the colimit \eqref{eq:15:germ-colimit} --- which shrinks the open unit disk
down to the origin --- makes it idempotent.
\end{remark}

This idempotent algebra shares the formal properties of $\mathbb{Z}[\![T]\!]$ and lets
us import Huber's theory of pairs into the gaseous context.

\begin{definition}[Analytification of an affine over a gaseous base]
\label{def:15:analytification}
Let $R$ be finite type over the underlying ring
$\mathbb{Z}[\widehat q^{\pm1}]^{\mathrm{gas}}(\ast)$, and equip it with the analytic
ring structure $(R, \operatorname{Mod}_R(\text{gaseous modules}))$. An element $f\in
R$ gives a section $\operatorname{Spec} R \xrightarrow{f} \mathbb{A}^1$. Moving the
germ-at-$0$ idempotent algebra of \cref{rmk:15:germs} to infinity gives a germ-at-$\infty$
idempotent algebra; one then passes to the subspace of $\operatorname{Spec} R$ on
which $f$ lands away from infinity --- concretely, where killing that idempotent
algebra after base change along $f$ holds,
\begin{equation}
\label{eq:15:analytic-condition}
(\text{germ-at-}\infty)\ \otimes_{R}\ R/f\ =\ 0 .
\end{equation}
The result is a new analytic space $\operatorname{Spec}(R)^{\mathrm{an}}$, the
\emph{analytification}.
\end{definition}

\begin{example}[Analytic $\Gm$]
\label{ex:15:gm-an}
For $R = \mathbb{Z}[T^{\pm1}]$ the analytification $\operatorname{Spec}(R)^{\mathrm{an}}
= \Gm^{\mathrm{an}}$ is exactly the object of \cref{lec:14} that one quotients by
$q^{\mathbb{Z}}$ to build the Tate curve (\cref{def:14:tate-construction}) --- up to
base change to a ring where $\mathbb{Z}$ is bounded, see \cref{rmk:15:analytic-corrections}.
\end{example}

\begin{remark}[Two corrections: which functions, and boundedness of the base]
\label{rmk:15:analytic-corrections}
As stated at the board, \cref{def:15:analytification} was corrected twice.
\begin{enumerate}
\item The analytification does \emph{not} depend on imposing
\eqref{eq:15:analytic-condition} for every $f\in R$ separately. It suffices to impose
it for the universal coordinate (for $\Gm$: for $T$ and $T^{-1}$); one automatically
gets the condition for all elements, by the general formalism. Imposing it na\"ively
for all $f$ is not even well-posed as a subset condition --- the locus ``away from
$\infty$'' is not closed under addition (the sum of two bounded functions need not be
bounded).
\item For the analytification of $\operatorname{Spec} R$ itself to be well-behaved one
must add the axiom that every element of the ground ring, in particular every integer
--- and $2$ being bounded is enough --- is bounded. Otherwise the analytification
changes the base: forcing all scalars to be bounded is itself a nontrivial operation.
With the axiom, $\operatorname{Spec} R$ is its own analytification. This axiom was
already present in the earlier write-up of complex-analytic geometry
\cite{clausen2020analytic}; the archimedean/Berkovich subtlety is exactly that the
archimedean branch of $\MBerk(\mathbb{Z})$ runs all the way out, so an unbounded
integer like $2$ escapes the ``away from $\infty$'' locus.
\end{enumerate}
\end{remark}

\begin{remark}[Solid and gaseous theories glue]
\label{rmk:15:glue}
The gaseous construction makes sense not only over
$\mathbb{Z}[\widehat q^{\pm1}]^{\mathrm{gas}}$ but over
$\mathbb{Z}[q]$ and $\mathbb{Z}[q]^{\mathrm{gas}}$: defining the structure only required
writing down the endomorphism $1 - q\cdot\mathrm{shift}$ of $P$ and asking it to
become an isomorphism, which needs neither $q$ topologically nilpotent nor $q$ a unit.
Specializing $q$ interpolates the states of matter (cf.\ \cref{rmk:12:states-of-matter}):
\begin{equation}
\label{eq:15:q-specialization}
q = 0\ \rightsquigarrow\ \text{uncompleted } \mathbb{Z}\text{-theory},\qquad
q = 1\ \rightsquigarrow\ \text{solid theory},\qquad
q\ \text{top.\ nilpotent unit}\ \rightsquigarrow\ \text{gaseous theory}.
\end{equation}
Working \emph{away} from the locus of topological nilpotence one forces both $q$ and
$q^{-1}$ to be gaseous, hence $1$ gaseous, hence solid; so over that locus
$\mathbb{Z}[q^{\pm1}]^{\mathrm{gas}}$ restricts to $\mathbb{Z}[q^{\pm1}]_{\solid}$.
Thus the solid and gaseous theories are glued together.
\end{remark}

\subsection{GAGA as an isomorphism of stacks}

\begin{remark}[A stack of analytifications]
\label{rmk:15:analytification-stack}
Any variable $q$ declared to be gaseous produces (via the germ-at-$0$ colimit
\eqref{eq:15:germ-colimit}, moved to $\infty$) an idempotent algebra and hence a
notion of analytification. Different choices of $q$ differing by a bounded unit ---
an element of $\Gm^{\mathrm{an}}$ --- give the \emph{same} algebra. So the parameter
space of notions of analytification is the stack
\begin{equation}
\label{eq:15:analytification-stack}
\operatorname{Spec}\bigl(\mathbb{Z}[\widehat q^{\pm1}]^{\,q\text{-}\mathrm{gas}}\bigr)\big/\Gm^{\mathrm{an}},
\end{equation}
a quotient of (the spectrum of) an analytic ring by an equivalence relation. This is a
new role of stacks in the theory: a stack parametrizing choices of analytic geometry
over a given base ring. For an analytic ring $R$, a map
\begin{equation}
\label{eq:15:gaseous-structure-map}
\operatorname{Spec} R\ \longrightarrow\ \operatorname{Spec}\bigl(\mathbb{Z}[\widehat q^{\pm1}]^{\,q\text{-}\mathrm{gas}}\bigr)\big/\Gm^{\mathrm{an}}
\end{equation}
is a \emph{gaseous structure on} $\operatorname{Spec} R$. (The quotient by
$\Gm^{\mathrm{an}}$ encodes that a topologically nilpotent unit should only matter up
to a bounded difference: given a Tate--Huber ring there exists such a unit, but it
plays no essential role beyond a bounded ambiguity. Clausen and Scholze noted at the
board that the group one divides by is a subtler object than the na\"ive
$\Gm^{\mathrm{an}}$ --- one modifies $q$ by units of norm $1$ or by raising to a power
--- and left the precise identification to a later lecture.
\todo{Precise group in \eqref{eq:15:analytification-stack}/\eqref{eq:15:gaseous-structure-map}:
$\Gm^{\mathrm{an}}$ vs.\ the subtler ``in-between'' group (units of norm $1$, or
modification of $q$ by raising to a power) --- the exact identification was
deferred by the lecturer to a later lecture.})
\end{remark}

\begin{remark}[Two functors and a natural transformation]
\label{rmk:15:two-functors}
Fix a gaseous structure \eqref{eq:15:gaseous-structure-map} on
$\operatorname{Spec} R$ such that $2\in\mathbb{Z}$ is bounded (\cref{rmk:15:analytic-corrections}).
Then one gets a theory of analytification: two functors from $R$-schemes to analytic
stacks over $R$, together with a natural transformation. For $X = \operatorname{Spec} A$
with $A$ an $R(\ast)$-algebra:
\begin{itemize}
\item on the one hand, view $X$ as an analytic stack over the uncompleted
$\mathbb{Z}$ and base change to $R$, giving $X_R := X\times_{\mathbb{Z}} R$ (ordinary
algebraic geometry over $R$);
\item on the other hand, form the analytification $X_R^{\mathrm{an}}$ over $R$
(requiring the coordinate functions to land away from $\infty$).
\end{itemize}
Both glue over non-affine $X$, and there is a natural map
\begin{equation}
\label{eq:15:gaga-nt}
\begin{tikzcd}[row sep=large]
X_R^{\mathrm{an}} \arrow[d, "\mathrm{nt}"] \\
X_R = X\times_{\mathbb{Z}} R .
\end{tikzcd}
\end{equation}
\end{remark}

\begin{theorem}[GAGA]
\label{thm:15:gaga}
With notation as in \cref{rmk:15:two-functors}, and assuming every $f\in R(\ast)$ is
bounded, if $X\to \operatorname{Spec}(R(\ast))$ is proper (and finitely presented),
then the natural transformation \eqref{eq:15:gaga-nt} is an isomorphism:
\begin{equation}
\label{eq:15:gaga-iso}
X_R^{\mathrm{an}}\ \xrightarrow{\ \sim\ }\ X_R .
\end{equation}
\end{theorem}

\begin{remark}[On the finite-presentation hypothesis]
\label{rmk:15:gaga-fp}
Clausen flagged the finite-presentation hypothesis as one he was unsure is needed for
the isomorphism itself: it entered when he last thought about setting up the formalism
with lower-shrieks and proper pushforward (to know, e.g., that resolving $R/I$
relatively over a noetherian $R$ gives compact projective generators that are products
of copies of $R/I$). It may be inessential for the GAGA statement proper but relevant
for proper pushforward ($f_* = f_!$); kept in to be safe.
\end{remark}

\begin{remark}[Formal and non-formal consequences]
\label{rmk:15:gaga-consequences}
\Cref{thm:15:gaga} implies \emph{formally} that the derived categories agree,
\begin{equation}
\label{eq:15:gaga-D}
D(X_R^{\mathrm{an}})\ =\ D(X_R),
\end{equation}
a version of GAGA; and \emph{non-formally} (under further, quite general technical
conditions on $R$) that, e.g.,
\begin{equation}
\label{eq:15:gaga-vect}
\operatorname{Vect}(X_R^{\mathrm{an}})\ =\ \operatorname{Vect}(X_R)
\ \Longrightarrow\ \text{classical GAGA.}
\end{equation}
\end{remark}

\begin{remark}[Special cases]
\label{rmk:15:gaga-special}
All classical GAGA theorems are special cases of \cref{thm:15:gaga}.
\begin{enumerate}
\item \emph{Complex-analytic GAGA} (Serre \cite{Serre_1956}): over
$\mathbb{C}^{\mathrm{gas}}$, or $\mathbb{C}$ with a $p$-liquid structure, etc.
\item \emph{Formal GAGA} (Grothendieck): for a complete noetherian ring with a proper
scheme over it, coherent sheaves agree with coherent sheaves on the formal completion
(compatible collections over the nilpotent thickenings, \eqref{eq:15:vect-formal}).
In Huber-pair terms one works over $\operatorname{Spa}(R^{\wedge}_I, R^{\wedge}_I)$,
which lives over $\Zsolid$ and inherits the closed-unit-disk notion of
analytification.
\item \emph{Non-archimedean GAGA} over $\mathbb{Q}_p$ (or any complete
non-archimedean field): take the analytic ring $\operatorname{Spec}(\mathbb{Q}_{p,\solid})$
with the gaseous structure corresponding to
\begin{equation}
\label{eq:15:nonarch-gaga}
\operatorname{Spec}(\mathbb{Q}_{p,\solid})\ \xrightarrow{\ q\mapsto p\ }\
\operatorname{Spec}\bigl(\mathbb{Z}[\widehat q^{\pm1}]^{\,q\text{-}\mathrm{gas}}\bigr),
\end{equation}
$p$ being the standard choice of topologically nilpotent unit. This picks out the
usual analytification of algebraic varieties over $\mathbb{Q}_p$.
\end{enumerate}
\end{remark}

\begin{remark}[The $\mathbb{Q}_p$ excursion, corrected]
\label{rmk:15:qp-overconvergent}
Clausen initially suggested that choosing instead a gaseous structure on
$\mathbb{Q}_{p,\solid}$ factoring through $\Zsolid$ would give a different GAGA
theorem, with the analytification of $\mathbb{A}^1$ becoming the closed unit disk.
This was corrected in discussion. With that ``$\Zsolid$'' structure one does \emph{not}
get a GAGA statement at all: the required boundedness axiom
(\cref{rmk:15:analytic-corrections}) would force $1/p$ to be bounded, and forcing all
scalars bounded in the solid $\mathbb{Q}_p$ setting kills everything (the $0$ ring).
(The false expectation ``$\mathbb{A}^{1,\mathrm{an}} = $ closed disk'' also fails
because forcing $T$ analytic gives no reason for $T/p$ to be analytic, contradicting
the $\Gm(\mathbb{Q}_p)$-rescaling symmetry.) The salvageable point: there \emph{is} a
second, genuine gaseous structure on $\mathbb{Q}_{p,\solid}$ obtained by factoring
through $\mathbb{Z}_{p,\solid}$, and the difference between it and the standard one
\eqref{eq:15:nonarch-gaga} is morally the difference between usual rigid geometry and
an over-convergent version of rigid geometry.
\end{remark}

\subsection{Relations between examples}

Beyond GAGA, one wants the category of analytic stacks to be large enough to house
\emph{maps} between the various classes of examples, whenever a comparison is
expected. Some maps to be made sense of as morphisms of analytic stacks:

\begin{remark}[Comparison maps]
\label{rmk:15:relations}
\begin{enumerate}
\item The Tate curve maps to the topological circle,
\begin{equation}
\label{eq:15:tate-to-circle}
\mathrm{Tate}_q\ \longrightarrow\ (0,\infty)/q^{\mathbb{Z}}\ =\ S^1,
\end{equation}
from Scholze's uniformization (\cref{def:14:tate-construction} and the square there),
now to be a map of analytic stacks: $S^1$ is an algebraic (indeed topological) stack
via \cref{ex:15:betti}, and algebraic stacks embed into analytic stacks even over the
initial analytic ring.
\item A complex-analytic space $X$ maps to its underlying topological space,
\begin{equation}
\label{eq:15:X-to-XC}
X\ \longrightarrow\ X(\mathbb{C}),
\end{equation}
which explains (essentially by definition) the well-known fact that the coherent
cohomology of $X$ localizes on the underlying topological space.
\item A Banach ring $R$ maps its (liquid) spectrum to the Berkovich spectrum,
\begin{equation}
\label{eq:15:banach-to-berkovich}
\operatorname{Spec}^{\mathrm{liq}}(R)\ \longrightarrow\ \MBerk(R),
\end{equation}
in the same way that the $(R,R^+)$ affine analytic stacks match Huber's picture.
\end{enumerate}
The comparison isomorphisms one expects between these (e.g.\ Betti vs.\ de Rham in the
complex-analytic setting) should likewise promote to isomorphisms of analytic stacks.
\end{remark}

\subsection{Set-theoretic technicalities}
\label{subsec:15:set-theory}

We at least pin down what the category of analytic \emph{stacks} is (the Grothendieck
topology is deferred).

\begin{remark}[The problem]
\label{rmk:15:set-theory-problem}
In the algebraic analogue, $\operatorname{Sh}^{\mathrm{fpqc}}(\operatorname{Aff})$ is a
full subcategory of $\operatorname{Fun}(\operatorname{CRing}, \mathrm{Set})$ (or
groupoids), but $\operatorname{Aff} = (\operatorname{CRing})^{\mathrm{op}}$ is not a
small category, so morphisms in the functor category can involve more than a set's
worth of data --- a genuine pain.
\end{remark}

\begin{definition}[Accessible presheaves]
\label{def:15:accessible}
Instead of $\operatorname{PSh}(\operatorname{Aff})$, consider the \emph{accessible}
presheaves
\begin{equation}
\label{eq:15:accessible-psh}
\operatorname{PSh}^{\mathrm{acc}}(\operatorname{Aff})\ =\ \operatorname{Fun}^{\mathrm{acc}}(\operatorname{CRing}, \mathrm{Set}),
\end{equation}
where a functor $F$ is \emph{accessible} iff there exists a regular cardinal $\kappa$
such that $F$ commutes with $\kappa$-filtered colimits.
\end{definition}

\begin{remark}[Reading the cardinal]
\label{rmk:15:reading-kappa}
One should think of a regular cardinal $\kappa$ not as a set but as indexing the sets
of smaller cardinality. For $\kappa = \aleph_0$ these are the finite sets and
$\kappa$-filtered colimits are ordinary filtered colimits (so a functor commuting with
filtered colimits is accessible). Larger $\kappa$ means \emph{fewer} $\kappa$-filtered
colimits, hence \emph{more} functors commute with them: e.g.\ $\kappa=\aleph_1$ indexes
countable sets, and a $\aleph_1$-filtered diagram is one in which every countable
subset admits a cone.
\end{remark}

\begin{remark}[Accessible $=$ small colimit of representables]
\label{rmk:15:accessible-representable}
For every $X\in\operatorname{Aff}$ the representable $h_X$ is accessible; in fact a
presheaf is accessible iff it is a \emph{small} colimit of representables $h_X$. Thus,
thinking of affine schemes as building blocks and colimits as gluing, accessibility
says precisely that one is only allowed to glue a set's worth of things at a time.
Moreover fpqc sheafification preserves accessible presheaves (Waterhouse
\cite{Waterhouse_1975}), so one may impose the sheaf condition without leaving the
accessible world and without set-theoretic obstruction.
\end{remark}

\begin{remark}[What makes this work]
\label{rmk:15:what-works}
Two non-trivial inputs:
\begin{enumerate}
\item $\operatorname{CRing}$ is presentable --- in fact compactly generated; this is what
makes the theory of accessible presheaves work (it is what lets ``commutes with
$\kappa$-filtered colimits'' match ``small colimit of representables'');
\item there is an a priori cardinality bound on fpqc covers: every fpqc cover in
$\operatorname{Aff}$ is an $\aleph_1$-filtered colimit of $\aleph_1$-compact
(countably presented) fpqc covers. Equivalently, every flat cover of a ring is an
$\aleph_1$-filtered colimit of countably-presented flat covers.
\end{enumerate}
Such a bound is what lets one verify that fpqc sheafification preserves accessibility.
\end{remark}

\begin{theorem}[Presentability of analytic rings]
\label{thm:15:anring-presentable}
The category of analytic rings is presentable; in fact it is $\kappa$-compactly
generated for $\kappa = (2^{\aleph_0})^+$ (so $\kappa$-small means of cardinality at
most the continuum).
\end{theorem}

This is what lets one formally avoid the set-theoretic difficulties of taking
presheaves and sheaves on the large category $(\operatorname{AnRing})^{\mathrm{op}}$.

\begin{remark}[An enrichment aside]
\label{rmk:15:enrichment}
The category of analytic rings is naturally enriched over condensed anima (its mapping
spaces are condensed anima). Analytic \emph{stacks} are, however, not enriched over
analytic rings (nor obviously over themselves). They \emph{do} carry a distinct
enrichment over light condensed sets: an $S$-valued point of a stack is a map
\begin{equation}
\label{eq:15:S-valued-point}
S\times X\ \longrightarrow\ Y,\qquad S\ \text{a light profinite (condensed) set},
\end{equation}
imported through the very trivial analytic rings of continuous functions
(\cref{ex:15:betti}). This enrichment has essentially no relation to the condensed-anima
enrichment on analytic rings, which Clausen notes is genuinely confusing --- for
instance Wies\l awa Nizio\l{} (in recent work on $p$-adic coefficient systems) wants
sheaves of solid abelian groups on the pro-\'etale site, so that profinite sets appear
in two orthogonal directions at once.
\end{remark}
